\section{Capítulo~\ref{sec:dishbrain}. DishBrain}

El diseño del banco y los resultados del juego provienen de un solo trabajo experimental y de su continuación; las fórmulas del ruido y de la deriva hacia los buenos ajustes se deducen en el capítulo y se comprueban con el modelo del libro, y el mecanismo del aprendizaje en la placa no está establecido.

{\footnotesize
\setlength{\tabcolsep}{3pt}\renewcommand{\arraystretch}{1.15}
\begin{longtable}{>{\raggedright}p{0.5cm}>{\raggedright}p{4.0cm}>{\raggedright}p{5.6cm}>{\raggedright}p{2.3cm}>{\raggedright\arraybackslash}p{1.7cm}}
\toprule
N.º & Afirmación & Experimento y qué se midió & Fuente & Estado \\
\midrule\endfirsthead
\toprule
N.º & Afirmación & Experimento y qué se midió & Fuente & Estado \\
\midrule\endhead
1 & Banco: unas $800\,000$ células de corteza de ratón o unas $10^6$ células humanas por chip; $26\,000$ electrodos en $8$~mm$^2$, registro de hasta $1024$, estimulación a través de $8$ & Métodos del trabajo: chip MaxOne, $26\,000$ electrodos de platino en $8$~mm$^2$, registro de $1024$ canales a $20\,000$ muestras por segundo; técnicamente se pueden estimular hasta $32$ electrodos, de forma independiente se logró con $8$; los cultivos de ratón se usaron desde unos $10$~días & \cite{Kagan2022} & medido \\
2 & Lazo con ciclo de $10$~ms: las espigas de las regiones motoras mueven la raqueta (fig.~\ref{fig:dishloop}) & Las espigas se cuentan en ventanas de $10$~ms; el retardo desde la espiga hasta el estímulo de respuesta es de unos $5$~ms y lo fija el búfer del equipo & \cite{Kagan2022} & medido \\
3 & Entrada: código de lugar (cuál de los $8$ electrodos) y de frecuencia, $4$--$40$~pulsos/s & En el experimento principal la frecuencia de estimulación crece linealmente de $4$~Hz, con la pelota junto a la pared lejana, a $40$~Hz junto a la raqueta; la amplitud de los pulsos de la pelota es de $75$~mV; los pulsos son rectangulares bifásicos & \cite{Kagan2022} & medido \\
4 & Salida $\Delta p=c\,(n_\uparrow-n_\downarrow)$~\eqref{eq:dishreadout}, ruido de la cuenta por ventana $1/\sqrt{RT}$ & La regla de la diferencia entre dos regiones está descrita en el trabajo (con pesos de las regiones según su actividad), pero no se da la fórmula exacta. La estimación del ruido es consecuencia del conteo de Poisson y se deduce en el capítulo & \cite{Kagan2022}; deducción en el capítulo & teoría \\
5 & Maestro: tras el acierto, un estímulo predecible de $75$~mV, $100$~pulsos/s, $0.1$~s; tras el fallo, uno impredecible, $150$~mV, $5$~pulsos/s, $4$~s en lugares e instantes aleatorios, y luego $4$~s de pausa & Métodos: tras el acierto, $75$~mV, $100$~Hz, $100$~ms en los $8$ electrodos a la vez; tras el fallo, $150$~mV, $5$~Hz en electrodos aleatorios e instantes aleatorios durante $4$~s, y luego $4$~s sin estimulación. En el cultivo no hay dopamina & \cite{Kagan2022} & medido \\
6 & La red actúa de modo que su entrada se vuelva predecible & El principio de energía libre es una propuesta teórica de la que los autores derivaron el protocolo; el experimento concuerda con él, pero no lo distingue de mecanismos más sencillos (filas~7--8) & \cite{Friston2010,Kagan2022} & hipótesis \\
7 & Si el fracaso sacude la red y el éxito no, el tiempo en un ajuste es $\rho\propto1/P_{\text{fallo}}$~\eqref{eq:dishdensity} (proposición~\ref{prop:dishscramble}) & Se sigue del balance de flujos en una cadena de Markov con empujones simétricos; se deduce en el capítulo. No hay comprobación directa en un cultivo & deducción en el capítulo & teoría \\
8 & El modelo de “revolver tras el fallo” aprende: fracción de aciertos $0.21\to0.38$; con empujón tras cada peloteo, $\approx0.12$; sin empujones, $0.19$ (fig.~\ref{fig:dishmodel}) & Cálculo, script \texttt{models/dishbrain\_model.py}, $40$ cultivos modelo con $2000$ peloteos cada uno: $0.21\to0.38$; $0.13\to0.11$; $0.19\to0.19$ & --- & modelo \\
9 & Las series de pelotas devueltas se alargan solo en cultivos vivos con el lazo completo; los controles no aprenden & $399$ sesiones de $20$~min: $9$ cultivos de ratón ($101$ sesiones), $11$ humanos ($138$), $6$ chips con medio ($80$), $20$ cultivos en reposo ($42$), $3$ simulaciones ($38$). La longitud media del peloteo en los últimos $15$~min es mayor que en los primeros $5$ solo en cultivos vivos; con células que no son neuronas (HEK293T) y con chips con medio no hay efecto & \cite{Kagan2022} & medido \\
10 & Un aumento significativo en la sesión, solo con realimentación completa; con silencio en lugar del estímulo, al final de la sesión se juega mejor que sin realimentación, pero el efecto es menor & $486$ sesiones en $3$ días: “estímulo” ($n=27$), “silencio” ($17$), “sin realimentación” ($15$). El aumento en el tiempo solo es significativo en la condición “estímulo”; al final de la sesión, “silencio” superaba a “sin realimentación” con un efecto menor & \cite{Kagan2022} & medido \\
11 & El efecto es pequeño; si la red aprende de forma duradera es una cuestión abierta & Dentro de la sesión la mejora es sólida, pero, según los propios autores, no se observó un aprendizaje estable entre sesiones a lo largo de varios días & \cite{Kagan2022} & medido \\
12 & Durante el juego la red se acerca al régimen crítico, y cuanto más cerca, mejor el juego & Análisis de avalanchas de actividad en los mismos cultivos: la entrada estructurada acerca la red al régimen crítico, y la calidad del juego se correlaciona con la cercanía a él; pero la criticidad sola, sin información sobre las consecuencias de las acciones, no basta para aprender & \cite{Habibollahi2023} & medido \\
13 & La “sintiencia” del cultivo no está demostrada & Artículo de respuesta de treinta investigadores: un cambio de comportamiento en lazo cerrado no demuestra ni inteligencia ni sensaciones; no hay un experimento que lo muestre & \cite{Balci2023} & hipótesis \\
14 & Cuarto control propuesto: un estímulo predecible tras el fallo, de la misma duración y energía (sección~\ref{sec:dishspec}) & Este experimento no se ha hecho; es una propuesta del libro & --- & hipótesis \\
\bottomrule
\end{longtable}}
